% 这是中国科学院大学计算机科学与技术专业《计算机组成原理（研讨课）》使用的实验报告 Latex 模板
% 本模板与 2024 年 2 月 Jun-xiong Ji 完成, 更改自由 Shing-Ho Lin 和 Jun-Xiong Ji 于 2022 年 9 月共同完成的基础物理实验模板
% 如有任何问题, 请联系: jijunxoing21@mails.ucas.ac.cn
% This is the LaTeX template for report of Experiment of Computer Organization and Design courses, based on its provided Word template. 
% This template is completed on Febrary 2024, based on the joint collabration of Shing-Ho Lin and Junxiong Ji in September 2022. 
% Adding numerous pictures and equations leads to unsatisfying experience in Word. Therefore LaTeX is better. 
% Feel free to contact me via: jijunxoing21@mails.ucas.ac.cn

\documentclass[11pt]{article}

/Users/ycp/Source/Courses/ca-lab/cdp_ede_local/report-template/preamble.tex

\begin{document}

%%%%%%%%%%%%%%%%%%%%%%%%%%%
%改这里可以修改实验报告表头的信息
\newcommand{\name}{}
\newcommand{\studentNum}{}
\newcommand{\major}{计算机科学与技术}
\newcommand{\labNum}{}
\newcommand{\labName}{}
%%%%%%%%%%%%%%%%%%%%%%%%%%%

\input{head.tex}

\section{}

\noindent
$\bullet$
\textbf{}

\section{实验过程中遇到的问题、对问题的思考过程及解决方法}


\noindent
$\bullet$
\textbf{}

\vspace{1ex}

\noindent
$\bullet$
\textbf{}

\section{对讲义中思考题（如有）的理解和回答}

\noindent
$\bullet$
\textbf{}


\vspace{1ex}

\section{实验所耗时间}

\newcommand{\labTime}{5}

在课后，花费了大约\underline{\makebox[5em][c]{\labTime}}小时完成此次实验。

\end{document}
